% ============================================================
%  Section 6 -- Conclusion
% ============================================================

\chapter{Conclusion}

\bibopsname{} repond a un besoin concret : evaluer de maniere rigoureuse des
architectures LLM pour le support IT Michelin. Le projet montre qu'une reponse
LLM directe ne suffit pas pour un cas industriel ou les decisions dependent de
sources internes, de contraintes de securite et de metriques operationnelles.
L'approche multi-agents, fondee sur ReAct, RAG, base de connaissances et statut
serveur, apporte un gain mesurable.

Sur le benchmark principal, le systeme multi-agents obtient un score composite
de \metric{75,2/100}, contre \metric{35,0/100} pour le LLM unique. Il reduit la
latence de \metric{63,7 \%}, divise le cout par \metric{7,91}, reduit les tokens
de \metric{88,3 \%} et conserve un score securite de \metric{10/10}. Ces
resultats valident l'interet de l'architecture outillee.

La conclusion reste cependant prudente : les deux architectures echouent au gate
de qualite, car le meilleur score moyen n'est que de \metric{3,8/10} pour un
seuil fixe a \metric{7/10}. Le projet ne prouve donc pas qu'un agent BibOps est
pret a remplacer un support IT humain. Il prouve plutot qu'un framework
d'evaluation peut identifier objectivement les progres, les faiblesses et les
risques avant de parler de deploiement.

La Racing Arena complete cette conclusion en montrant que la securite des agents
est une question d'architecture. Dans le scenario observe, l'equipe validee
n'execute aucune injection et ne fuite aucune strategie, alors qu'une equipe
ReAct exposee execute 93 \% des injections recues. Ce resultat justifie
l'integration des tests adversariaux dans le cycle de benchmark.

Le travail accompli fournit donc une base solide :

\begin{itemize}
  \item un agent IT traceable et outille ;
  \item une evaluation multi-criteres explicite ;
  \item des artefacts JSON et graphiques exploitables ;
  \item une suite de 552 tests unitaires passes ;
  \item une simulation adversariale pour analyser les risques agentiques.
\end{itemize}

Les prochaines etapes sont claires : augmenter le corpus et le nombre de tickets
evalues, ameliorer la qualite des prompts et du RAG, calibrer les jugements avec
des experts humains, durcir la securite des endpoints et rendre la generation du
rapport entierement reproductible. Avec ces ameliorations, \bibopsname{} peut
devenir un outil de decision robuste pour comparer et gouverner des assistants
LLM en environnement industriel.
